\section{Coverage ledger and the exact boundary of the result}
\label{sec:coverage}

The requested Wikipedia and MathWorld articles~\cite{wiki,mathworld}
contain definitions, analytic identities, numerical tables, historical
remarks, implementation discussions, and deep arithmetic assertions.
The following ledger records \emph{62 result families developed here}.
Some entries contain several related equations; the number is a coverage
index, not a count of independent new theorems. The standard formulas
were cross-checked against the corresponding NIST chapter~\cite{dlmf}.
The contribution here is the finite constructive argument and its
arithmetic interpretation, not novelty of the gamma identities.

All entries have the fixed-name, finite-precision meaning explained in
\cref{sec:foundations}. ``Complex'' includes explicit separation from
every denominator's zero set. ``Positive'' includes a supplied positive
rational lower bound. Individual derivative and asymptotic orders can
be read as externally fixed, as explained in
\cref{sec:formalization}. The human-readable proofs are present; the
full Lean object-theory formalization remains the specified work.

\small
\begin{longtable}{@{}r >{\raggedright\arraybackslash}p{0.52\textwidth}>{\raggedright\arraybackslash}p{0.34\textwidth}@{}}
\toprule
ID & Result family & Proof / domain qualification\\\midrule\endhead
1 & Euler's constant and its harmonic-log error bound & (C1)--(C2), error $1/N$.\\
2 & Euler's limit for gamma & (C15), bounded pole-separated inputs.\\
3 & Weierstrass product for $1/\Gamma$ & (C7)--(C11), every disk.\\
4 & Entire analyticity of the reciprocal & (C16)--(C19), coefficient majorants.\\
5 & Reciprocal zeros exactly $0,-1,-2,\ldots$, all simple & (C22)--(C23) and finite root location.\\
6 & Gamma has no zeros & Positive lower bound $1/M_R$.\\
7 & Meromorphic gamma; simple poles and residues $(-1)^n/n!$ & (C24)--(C25), punctured neighborhoods.\\
8 & Functional equation and normalization & (C20)--(C21).\\
9 & Integer factorial values & (C21), finite recurrence.\\
10 & Complex conjugation and real values on the real domain & (C26).\\
11 & Euler's integral and absolute convergence & \cref{sec:integrals}, explicit endpoint tails.\\
12 & Positivity and computable upper and lower bounds & \eqref{eq:gamma-positive}--\eqref{eq:gamma-real-bound}.\\
13 & $|\Gamma(\sigma+iy)|\le\Gamma(\sigma)$ & Integral triangle inequality, $\sigma>0$.\\
14 & All derivatives as integrals with logarithmic factors & \cref{sec:integrals}, fixed derivative order.\\
15 & Exponential, unit-interval, and quadratic integral forms & \eqref{eq:exponential-substitution}--\eqref{eq:quadratic-substitution}.\\
16 & General positive scaling integral & \eqref{eq:general-scaling}, $c,r>0$.\\
17 & Lower and upper incomplete gamma: sum, recurrence and integer formula & End of \cref{sec:integrals}.\\
18 & Euler reflection and $\Gamma(z)\Gamma(-z)$ & (R1)--(R3), explicit sine margin.\\
19 & Gauss multiplication & (M1)--(M5), arbitrary positive integer multiplier.\\
20 & Legendre duplication & (M6), independent of Stirling.\\
21 & Triplication & (M7).\\
22 & Beta--gamma relation & (B1)--(B8), both real parts positive.\\
23 & Beta symmetry and recurrence & (B2)--(B5).\\
24 & Beta integral evaluations and trigonometric forms & (B9)--(B13), stated half-plane margins.\\
25 & $\Gamma(1/2)=\sqrt\pi$ & (R4), positivity selects the root.\\
26 & All positive and negative half-integer values & (V1).\\
27 & Third, quarter, fifth, sixth and twelfth relations & (V2)--(V6), positive roots.\\
28 & Product of all fractional arguments $r/m$ & (V7).\\
29 & Rising factorial and generalized binomial identities & (V8)--(V9), polynomial cancellation at apparent poles.\\
30 & Imaginary and half-shifted imaginary modulus formulas & (VM1)--(VM5), with a zero-safe version.\\
31 & Digamma partial fractions and their error & (C27)--(C28).\\
32 & All polygamma sums, recurrences, and sign patterns & (C29)--(C35).\\
33 & Integer and half-integer derivative values & (C32)--(C34), $\Gamma'(1)=-\gamma$.\\
34 & Differentiated reflection and multiplication & (D1)--(D3).\\
35 & Digamma values at thirds, quarters and sixths & (D4).\\
36 & Gamma derivatives by finite polynomial recursion & (C36).\\
37 & Taylor expansion at one and local log-gamma series & (C37)--(C41), disk radius $<1$.\\
38 & Entire reciprocal Maclaurin series and coefficient recurrence & (C19), (C42).\\
39 & Laurent constant and quantitative remainder at every pole & (C43).\\
40 & Expansions about every positive integer & (C44)--(C45), disk radius $<n$.\\
41 & Vertical-line modulus infinite product & (C46), real part separated from poles.\\
42 & Balanced finite-parameter infinite products & (C47)--(C48).\\
43 & Strict log-convexity and strict convexity & \eqref{eq:strong-log-convexity}, quantitative gaps.\\
44 & Effective Bohr--Mollerup uniqueness & \eqref{eq:bm-width-general}, fixed function presentations.\\
45 & Wendel and Gautschi inequalities & \eqref{eq:wendel-constructive}--\eqref{eq:gautschi-constructive}.\\
46 & Gamma-ratio asymptotics with a rate & \eqref{eq:general-gamma-ratio}, real shifts.\\
47 & Unique positive minimum and a PR algorithm for its location & \cref{sec:gamma-critical-points}.\\
48 & Sign alternation and one extremum between every two negative poles & \cref{sec:gamma-critical-points}.\\
49 & Stirling leading term and explicit relative error & \eqref{eq:stirling-first-bound}--\eqref{eq:stirling-multiplicative-bound}, $x>0$.\\
50 & Robbins factorial inequalities, and a real-variable strengthening & \eqref{eq:robbins-real}, $x\ge1$.\\
51 & Every order of the log-Stirling expansion with sharp signed remainder & \eqref{eq:sharp-stirling-remainder}, fixed order, real $x>0$.\\
52 & Every order of the digamma asymptotic expansion & \eqref{eq:sharp-digamma-remainder}, fixed order, real $x>0$.\\
53 & Multiplicative Stirling coefficients and an error-producing algorithm & End of \cref{sec:all-order-stirling}.\\
54 & Canonical analytic log-gamma on the right half-plane & \eqref{eq:canonical-loggamma}; branch specified.\\
55 & Binet's first and second formulas & \cref{sec:binet}, complex $\Rea z>0$.\\
56 & Complete monotonicity of the real Stirling remainder & Positive Laplace kernel in \cref{sec:binet}.\\
57 & Raabe and the unit-interval log-gamma integrals & \eqref{eq:unit-loggamma}--\eqref{eq:raabe}.\\
58 & Gamma--zeta Mellin identity & \eqref{eq:gamma-zeta-mellin}, $\Rea s>1$.\\
59 & Wallis product and a rational convergence rate & \cref{sec:further}.\\
60 & Gaussian moments and gamma-density normalization and moments & \cref{sec:further}.\\
61 & The two gamma-crossing equations and the golden ratio solution & \cref{sec:further}, valid gamma domains.\\
62 & Explicit alternative factorial interpolations satisfying recurrence & $\Gamma(z)e^{c\sin(2\pi z)}$; \cref{sec:further}.\\
\bottomrule
\end{longtable}
\normalsize

\subsection{Statements deliberately not claimed}

The scope boundary is mathematical, not a claim that the omitted
statements are impossible in iPRA. We do not provide proofs here of
H\"older's differential-transcendence theorem; irrationality,
transcendence, or algebraic-independence results for special gamma values;
elliptic-integral singular-value evaluations; Kummer's Fourier series;
the displayed continued fractions; the correctness and performance of
particular Lanczos or Spouge implementations; or the zeta functional
equation. The complex-sector form of the full Stirling expansion is
also outside the proved scope: the sharp all-order remainder theorem
here concerns the positive real axis. Binet's formulas do hold on the
certified complex right half-plane as explicitly proved.

Hankel contour representations are not included: proving them would
require an additional finite contour development with branches and
contour-error estimates. Historical statements, specific software
implementations, and long decimal tables are not mathematical theorem
targets of the present argument. The exact certificate run illustrates
the evaluator rather than certifying the source articles' decimal
tables or their quoted continued-fraction digits.

Finally, the first-order claim is the explicit arithmetic theorem schema
for certified primitive recursive presentations. It does not quantify
internally over all classical complex numbers or all real functions.
Nor does the archive claim a completed Lean derivation for all the
listed analysis. These distinctions are part of the statement of the
result, and preserve the constrained constructive setting requested.
